\chapter{转录：从 DNA 制作可使用的副本}

\begin{kaichang}
假设你家有一本祖传菜谱，牛皮封面，纸页发黄，上面还有太奶奶手写的批注。全家立下规矩：这本书\textbf{绝对不能带进厨房}——厨房里有油、有水、有火，一旦弄脏弄破，世上再没有第二本。可是今天轮到你做那道红烧肉，你需要其中一页。怎么办？

答案你早就知道：拿一张普通便签纸，把那一页\textbf{抄下来}，带进厨房。做完菜，纸条弄脏了也不要紧，扔掉就是；原书安安稳稳躺在书柜里，下次需要还能再抄一份。

先别急着说这太简单。把场景搬进你身体的任何一个细胞：细胞的“祖传菜谱”是 DNA，整套只此一份，珍贵到你全身每个细胞都要为它配一个带膜的“保险柜”——细胞核（第 66 章讲过它的结构为什么适合当档案）；而真正干活的“厨房”，也就是制造蛋白质的核糖体，全都分布在细胞核\textbf{外面}的细胞质里。菜谱不能出库，厨房却在库外，这道“抄写”工序是怎么完成的？用什么纸、什么字母？谁来抄？抄错了怎么办？

先猜一个答案，写在纸上。这一章结束时，你会看到细胞的做法和你抄菜谱惊人地像——只不过它的“抄写员”是一台蛋白质机器，抄出来的纸条叫 RNA。
\end{kaichang}

\section{两个房间之间的搬运问题}

老规矩，先只做观察，不急着解释。

二十世纪五十年代的生物学家手里有几条硬事实。第一条：用染料把细胞染上色，DNA 几乎全部待在细胞核里，几乎不出核门。第二条：制造蛋白质的场所——核糖体——在细胞质里，核糖体上热火朝天合成蛋白质的时候，附近并没有 DNA。第三条：细胞里除了 DNA，还有另一种长得很像的核酸，叫 RNA（核糖核酸）。它在细胞核里出现，也在细胞质里出现，\textbf{两头都露面}。第四条：一个细胞合成蛋白质越卖力，它里面的 RNA 就越多。拼命生产消化酶的胰腺细胞、快速分裂的细菌，RNA 含量都高得出奇；而那些几乎不造蛋白质的休眠细胞，RNA 也少得可怜。

还有一条更妙的观察。科学家让细胞短暂地“吃”一会儿带放射性的尿嘧啶——一种 RNA 特有的原料（下面马上讲它特殊在哪），然后追踪放射性的位置：几分钟后，放射性先出现在细胞核里，过一阵子才慢慢出现在细胞质的核糖体上。像是有什么东西，在核里被制造出来，然后被搬运到了厨房门口。

把这几条拼起来，一个猜想的轮廓已经很清楚了：DNA 不出核，却要把信息送到核外的核糖体；RNA 恰好两头跑，含量又和蛋白质生产的忙碌程度同步。信息很可能是这样走的：DNA 留在核里当底本，细胞照着需要的某一页，做出一份 RNA 副本，副本出核，送到核糖体。这份“可以带出库的工作副本”，就是这一章的主角。这个从 DNA 抄写出 RNA 的过程，叫\textbf{转录}。

\section{RNA：一个“轻装版”的 DNA}

把镜头推进到粒子层。RNA 和 DNA 有多像，又差在哪？第 66 章讲过，DNA 是一条由核苷酸串成的长链：每个核苷酸 = 一个磷酸 + 一个五碳糖 + 一个含氮碱基，糖和磷酸交替连成骨架，碱基伸在外面当“字母”。RNA 的基本构造一模一样，只有三处改动，但每一处都有后果。

\begin{itemize}[itemsep=2pt]
\item \textbf{糖多带了一个氧}。DNA 骨架里的糖叫脱氧核糖，RNA 里的叫核糖——差别就在糖环的第二个碳上：核糖那里多挂了一个羟基（\chem{-OH}）。多出来的这个羟基是个“活跃分子”，让 RNA 链更容易被化学攻击、被酶拆解。换句话说，RNA \textbf{天生短命}。
\item \textbf{字母 U 替换了 T}。DNA 的四个字母是 A、T、G、C；RNA 用 A、U、G、C。U（尿嘧啶）的结构和 T（胸腺嘧啶）几乎一样，照样和 A 配对——你可以把 U 理解成“简装版”的 T，省掉了一个甲基小尾巴。
\item \textbf{通常只有一条链}。DNA 是两条链拧成的双螺旋，RNA 大多数时候单链行动。单链没有“另一半”保护，更没有备份；但它也因此能自由折叠，团成各种三维形状——这一点后面会变成 RNA 的超能力。
\end{itemize}

\begin{table}[htbp]
\centering
\begin{tabular}{lll}
\toprule
 & DNA & RNA \\
\midrule
五碳糖 & 脱氧核糖（少一个羟基） & 核糖（多一个羟基） \\
碱基字母 & A、T、G、C & A、U、G、C \\
链的常态 & 双链双螺旋 & 多为单链，可折叠 \\
化学稳定性 & 高，适合长期存档 & 低，适合短期使用 \\
在细胞里的角色 & 底本、档案 & 工作副本、工具、零件 \\
\bottomrule
\end{tabular}
\end{table}

现在回头看这三处改动，会发现它们像商量好了一样，全都指向同一个用途：RNA 是一种\textbf{消耗品}。底本要稳定，要双链互相备份，所以用 DNA；工作副本量大、用完就扔、需要随做随销，所以用短命又便宜的 RNA。这不是细胞“想”出来的分工，而是自然选择留下的账本：长期存档的信息若保存在容易损坏的分子上，错误会累积到致命的程度；临时副本若用昂贵稳定的材料大量制造，又是纯粹的浪费。稳定性上的差异，对应的是使用方式上的差异。

\section{抄写员的工作手册}

规则层来了。负责抄写的机器叫 \textbf{RNA 聚合酶}——“聚合酶”的意思是“把小零件聚合成链的酶”（第 49 章讲过，酶是改变反应路径、降低启动能垒的蛋白质机器，它不提供能量）。它的工作手册可以写成五条。

第一条，\textbf{不是从书的第一页抄起，而是哪里写着“从这开始”就抄哪里}。DNA 上每个基因的开头不远处，都有一小段特殊的碱基序列，叫\textbf{启动子}，作用相当于书页上的“从此处抄起”标签。RNA 聚合酶在 DNA 上滑行，认出启动子就牢牢结合上去。这意味着转录是\textbf{选择性}的：细胞此刻需要哪种蛋白质，就只抄对应的那几个基因，而不是把整套基因组从头到尾复印一遍——那是复制（第 67 章）干的事，两道工序千万别混。

第二条，\textbf{只以一条链为模板}。聚合酶把结合处的双螺旋局部“撬开”一个小口，形成十几对碱基宽的“转录泡”，然后只读其中一条链（叫模板链），按碱基互补的老规矩配字母：模板上的 A 配 U，T 配 A，G 配 C，C 配 G。

第三条，\textbf{抄写的方向永远是 $5'$ 端到 $3'$ 端}。第 66 章讲过 DNA 链有方向性，RNA 也一样：新链只能从 $5'$ 端向 $3'$ 端延长，聚合酶一边读模板一边往前挪，身后被撬开的双螺旋重新合拢，新抄好的 RNA 单链从泡里“吐”出来，越拖越长。

第四条，\textbf{能量从原料里来}。聚合的零件是四种核苷三磷酸（NTP）——你可以把它们理解成“自带电池的字母”：每个零件接入链时，身上两个磷酸组成的尾巴脱落，整个过程让体系总能量下降，反应因此能自己向前走（和第 50 章 ATP 的道理一样：能量不藏在某一根键里，藏在整个反应的前后差里）。所以转录不需要外加能量，原料本身推着链往前长。

第五条，\textbf{抄到“完”字就停}。基因末尾有一段终止信号。在细菌里，常见的一种设计很巧妙：终止序列让刚抄出来的 RNA 自己折叠成一个“发卡”加一串 U——RNA 单链能折叠的超能力在这里派上了用场——这个结构会把聚合酶“别停”，让它松手，RNA 释放出来。一次转录到此结束。

最后补一条容易被忽略、却极重要的规则：\textbf{同一个基因可以反复抄}。一份 RNA 副本在细菌里往往几分钟就被拆掉了，可只要细胞还需要这种蛋白质，聚合酶就可以一次又一次回到那个启动子，抄出几十、几百份副本。底本岿然不动，副本源源不断。

\begin{qiaoliang}
感受一下“选择性抄写”省了多少事。大肠杆菌的基因组约有 $4.6\times 10^{6}$ 对碱基，整个复制一遍（第 67 章）就算以每秒约 1000 个碱基的高速，也要四十分钟上下。而一个普通基因的编码区大约 1000 个核苷酸，细菌的 RNA 聚合酶每秒抄 40～80 个字母——\textbf{抄完一个基因只要十几到二十几秒}。再算一笔产出账：细菌里一份 mRNA 的寿命大约几分钟，在被拆解前，它可以被核糖体反复读取，生产出几十份相同的蛋白质。也就是说：抄一页纸二十秒，换来几十份蛋白质产品；要是每次都复印整套 460 万字母的“全书”，光是抄写就要四十分钟，还制造出几百万个此刻根本用不上的字母。选择性，是这套信息系统里最值钱的规则。
\end{qiaoliang}

\section{把抄写规则写成符号}

现在才轮到符号登场。假如一小段 DNA 模板链的序列是

\[ \text{模板链（DNA）：}\quad \text{3'-TAC GGA CTA-5'} \]

按“A 配 U、T 配 A、G 配 C、C 配 G，方向相反”的规则，抄出来的 RNA 是

\[ \text{信使 RNA：}\quad \text{5'-AUG CCU GAU-3'} \]

注意三件事：RNA 里没有 T，模板上的 A 变成了 U；两条链方向相反（模板从 $3'$ 读到 $5'$，新链从 $5'$ 写到 $3'$）；抄出来的这条 RNA 序列，和 DNA 里\textbf{另一条}没当模板的链几乎一模一样，只是 T 换成了 U——所以那条链常被叫作“编码链”。

下面这幅图把整个过程画在一起。老规矩先声明：这是人为的示意图，圆圈和线的大小、形状都不代表分子的真实长相。

\begin{center}
\begin{tikzpicture}[scale=1,every node/.style={font=\small}]
\draw[thick] (0,0.35) -- (8,0.35);
\draw[thick] (0,-0.35) -- (8,-0.35);
\draw[thick] (3,0.35) .. controls (3.6,1.0) and (4.4,1.0) .. (5,0.35);
\draw[thick] (3,-0.35) .. controls (3.6,-1.0) and (4.4,-1.0) .. (5,-0.35);
\draw[fill=gray!25] (4,0) circle (0.6);
\node at (4,0) {聚合酶};
\draw[thick,->] (4,0.62) .. controls (4.5,1.3) .. (5.7,1.5);
\node at (6.9,1.5) {新生的 RNA};
\node at (1.2,0.8) {DNA 双链};
\draw[->] (7.0,0.75) -- (8.2,0.75);
\node at (7.6,1.1) {移动方向};
\end{tikzpicture}
\end{center}

聚合酶骑在 DNA 上，把双链撬开一个泡，RNA 从泡里长出来，身后的双链重新合上。

\begin{shiyan}
\textbf{用纸条扮演一次 RNA 聚合酶}（纸笔游戏，安全）

材料：一张白纸、一支笔、一把剪刀、一把尺子。

步骤：①在纸上画一条 12 格的纸条，在格子里从 $3'$ 端向 $5'$ 端写一段“DNA 模板链”，例如 TAC GGA CTA TGC。②拿另一张纸条，从 $5'$ 端向 $3'$ 端，按配对规则逐格“抄写”出 RNA：见到 A 写 U，见到 T 写 A，见到 G 写 C，见到 C 写 G。③把两条纸条并排摆好，检查方向是不是相反的、RNA 里是不是一个 T 都没有。④故意抄错一格（比如把某个 C 写成 A），再看看后面整行会发生什么。

观察点：你“读”模板的方向和“写”新链的方向，为什么必须相反？如果忘了把 T 换成 U，这条“RNA”和 DNA 还有什么区别？把抄错的纸条留着，第 72 章讲突变时它会派上用场。

安全提示：使用剪刀请小心。这只是一个帮助思考的模型——真实的聚合酶每秒抄几十个字母，而且不会“看错格子看半天”。
\end{shiyan}

\section{“一个基因一种酶”：一个伟大的简化}

转录规则讲完了，现在该面对本章最容易被讲歪的一句话。

1941 年，美国遗传学家比德尔和塔特姆用红色面包霉（一种真菌）做实验：他们用 X 射线轰击孢子制造突变，发现有些突变菌株失去了合成某种营养物质的能力，而且每破坏一个基因，恰好只坏掉合成路线上的一个酶。于是他们提出：\textbf{一个基因负责制造一种酶}。在那个连基因的化学本质都还没查明的年代，这是划时代的进步——它第一次把“基因”这个抽象的遗传因子，和“蛋白质”这个具体的分子机器连在了一起。生命这台化学城市里，图纸和机器之间终于有了桥。

但科学史的诚实之处在于：这句话后来一修再修。第一修，改成“一个基因对应一条多肽链”——因为人们发现，像血红蛋白这样的蛋白质由两种不同的链拼成，两种链各自由不同基因制造；第二修，发现有的基因产物根本不是蛋白质，而是 RNA 本身；第三修，发现真核生物的一个基因可以经过不同“裁剪”，产生好几种不同的蛋白质（马上细讲）。到今天，基因和蛋白质的关系更接近一张\textbf{多对多的网}：一个基因可能产生多种产物，一种蛋白质可能需要多个基因的产物拼装。

这不是说比德尔和塔特姆错了。在他们的时代、在他们研究的那些酶上，一一对应相当常见，这个简化让遗传学第一次变得可以做实验、可以算数。所有教科书里的模型都是这样：先给你一个能抓住主要矛盾的简化版，等你站稳了，再告诉你边界在哪里。这一章剩下的部分，就是带你去看边界。

\section{那份“信使”是怎样被找到的}

\begin{zhengju}
二十世纪五十年代末，生物学卡在一个矛盾上。一方面，所有证据都指向“蛋白质在核糖体上制造”；另一方面，核糖体里的 RNA（核糖体 RNA）在所有细胞里都长得差不多——如果核糖体本身就是图纸，那它怎么能今天造胰岛素、明天造血红蛋白？当时有两种针锋相对的猜想：其一，每种蛋白质都有自己专属的核糖体，图纸就焊在核糖体上；其二，核糖体只是一台通用机器，图纸以某种短命的“信使”形式临时插进来，用完即走。

1961 年，布伦纳、雅各布和梅塞尔森用一个精巧的实验一锤定音。他们让 T2 噬菌体去感染大肠杆菌（就是第 65 章赫尔希—蔡斯实验里的那种病毒）：噬菌体一进细胞，细菌就停工转产病毒蛋白。关键问题是：病毒蛋白是在\textbf{旧的}细菌核糖体上造的，还是在\textbf{新的}病毒核糖体上造的？

他们的办法是给核糖体“称体重”。让细菌先在含重同位素（$^{15}\mathrm{N}$、$^{13}\mathrm{C}$）的培养基里长大，全身核糖体都是“重”的；感染时把细菌转到普通轻培养基里，同时加入放射性标记，给新合成的分子打上记号。如果新造了病毒核糖体，它应该是“轻”的、带放射性；如果病毒只是往旧核糖体里插了信使，那么放射性应该出现在“重”核糖体上。结果清清楚楚：带放射性的、感染后才出现的病毒 RNA，全部骑在\textbf{旧的、重的}细菌核糖体上。新核糖体一个都没造。

结论：核糖体是通用机器，真正换班的是一份份临时插进来的 RNA 图纸。同年，雅各布和莫诺为它起名——\textbf{信使 RNA}（mRNA）。回看历史，其实更早就有线索：沃爾金等人早就观察到感染后会出现一种更新极快、比例恰好反映病毒 DNA 的 RNA，只是当时没人敢断言它就是图纸。科学的证据链常常这样：线索先到，能一锤定音的实验设计后到，而“敢把两份证据接起来”的解释最后才到。
\end{zhengju}

\section{RNA 家族：信使只是其中一员}

mRNA 出尽了风头，但细胞里的 RNA 远不止这一种。粗略分一下工：

\begin{itemize}[itemsep=2pt]
\item \textbf{信使 RNA（mRNA）}：从 DNA 抄来的工作图纸，携带制造蛋白质的序列信息。
\item \textbf{核糖体 RNA（rRNA）}：核糖体的主体。前面说“制造蛋白质的是核糖体”，更精确地说，核糖体里真正催化氨基酸成键的活性中心，主要是 rRNA，蛋白质反而多是辅助支架。也就是说，\textbf{最核心的蛋白质制造机，本身是一个 RNA 机器}。这个发现让“RNA 只是便条纸”的印象彻底破产，也是“生命早期可能由 RNA 既当档案又当工具”这一假说（第 56 章提过）的重要证据之一。
\item \textbf{转运 RNA（tRNA）}：一种折叠成特定形状的小 RNA，一头认 mRNA 上的字母，一头抓着对应的氨基酸，是“核酸语言”和“氨基酸语言”之间的翻译适配器。它怎样工作、四种字母怎样变成二十种氨基酸，是第 69 章的主角，这里只记住它的存在。
\item \textbf{调控 RNA}：一大类不制造蛋白质、专管“开关和音量”的 RNA。比如 microRNA，只有二十来个字母长，能配对结合到特定的 mRNA 上，让它翻译不出来或者提前销毁——相当于给便条盖“作废”章。这类 RNA 是基因调控的重要角色，第 70 章会细讲细胞怎样决定什么时候转录、转录多少。
\end{itemize}

看明白了吗：从 DNA 转录出来的产物，并不都是为了变成蛋白质。有的 RNA 本身就是最终产品——是机器、是零件、是开关。这已经是“一个基因对应一种蛋白质”这句老话的第一道裂缝。

\section{真核细胞的“后期制作”}

到目前为止，我们的故事主要在细菌的舞台上讲——细菌没有细胞核，DNA 直接泡在细胞质里，RNA 一抄出来马上就能被核糖体读，甚至抄到一半翻译就开工了，转录和翻译在同一条流水线上。

你、我、院子里的猫、窗台上的绿萝是真核生物，规矩要多得多。我们的 DNA 锁在核里，刚抄出来的 RNA 还不能直接用，要先经过三道“后期制作”，才能拿到出核通行证。

\textbf{第一道，戴帽子。}新 RNA 的 $5'$ 端一冒头，就被加上一个经过特殊修饰的鸟嘌呤“帽子”。这顶帽子有两个作用：像安全帽一样保护 RNA 前端不被酶啃掉，又像工牌一样被核糖体识别——没戴帽子的 RNA，核糖体不认。

\textbf{第二道，加尾巴。}RNA 的 $3'$ 端被接上一长串腺嘌呤，叫多聚 A 尾，几十到几百个 A 连成一片。尾巴像一根“定时保险丝”：拆解酶从尾巴开始啃，尾巴越长，这份 mRNA 活得越久，能被翻译的次数越多。同一种 mRNA，尾巴长短不同，产能就不同——调控的暗门之一。

\textbf{第三道，也是最惊人的一道：剪接。}1977 年，科学家在研究一种感染人细胞的病毒时，撞上了一件怪事：病毒的 mRNA 比它的 DNA 模板\textbf{短}，怎么对都对不上。最后真相大白：真核生物的基因大多是“打断的”——有用的段落（叫\textbf{外显子}）之间，隔着大段不进入最终 mRNA 的序列（叫\textbf{内含子}）。转录先把整段照抄下来，然后由一台叫“剪接体”的机器（它的核心零件又是 RNA）把内含子整段剪掉，把外显子首尾拼接起来。这项发现震惊了当时的生物学界，夏普和罗伯茨因此在 1993 年获得诺贝尔奖。

而剪接最妙的地方在于：\textbf{同一份原始转录本，可以剪出不同的成品}。哪些外显子保留、哪些跳过，是可以变化的——这叫\textbf{可变剪接}。一个基因，在肌肉细胞里这样剪，在神经细胞里那样剪，就产生两种氨基酸序列不同、功能略有差异的蛋白质。极端的例子是果蝇一个叫 Dscam 的基因，靠可变剪接理论上能产生三万八千多种不同的 mRNA——比果蝇全部基因的总数还多。人类只有约两万个蛋白质编码基因，体内的蛋白质种类却远超两万，可变剪接是重要的扩容手段之一。

\begin{wuqu}
“一个基因对应一种蛋白质。”——这句流传了半个多世纪的话，是历史上的简化，不是今天的结论。反例随手就有：①果蝇 Dscam 基因靠可变剪接能产生上万种不同的蛋白质版本；②血红蛋白由两种不同的多肽链拼成，分别来自两个基因——是“多个基因一种蛋白质”；③细胞里大量基因的产物是 rRNA、tRNA、microRNA 这类功能性 RNA，一辈子不变成蛋白质。更准确的说法是：基因是可以产生功能产物的 DNA 区段，产物和蛋白质之间是多对多的关系。那为什么教科书还常说那句老话？因为在细菌里、在许多经典研究的案例里，它大体成立，是个很好用的入门模型——但请把它当作脚手架，不是大楼本身。第 71 章会把“基因到底是什么”这个问题彻底拆开重审。
\end{wuqu}

\section{这套模型的边界}

每个模型都有保质期，本章的“抄写模型”也不例外，有三处需要你保持警觉。

第一，我们把转录画成一台机器沿着轨道匀速前进，真实画面要乱得多：聚合酶会卡住、倒退、被别的蛋白质催着走或拦下来；启动子有强有弱，同一个基因在不同细胞、不同时刻的抄写频率天差地别。什么时候抄、抄多快，是被一整套调控系统实时决定的——那是第 70 章的主题。

第二，“DNA 的信息流向 RNA，再流向蛋白质”这个方向，在绝大多数细胞里成立，但不是宇宙铁律。HIV 这类病毒带着一种叫逆转录酶的机器，能反过来以 RNA 为模板抄写 DNA，把自己的基因插进宿主的基因组——箭头被整个反了过来。流感、新冠这类病毒干脆用 RNA 当基因组，遗传信息根本不经 DNA 的手。中心法则的“箭头图”是主干，不是单行线。

第三，我们默认“抄出来的副本内容总是忠于底本”。其实 RNA 聚合酶的出错率大约是每抄一万到十万个字母错一个，比 DNA 聚合酶（第 67 章讲过的带校对功能的抄写员）高出上百倍，而且转录没有复制那样严密的校对系统。为什么细胞受得了？因为副本是消耗品：错一份 RNA，顶多生产出一份残次蛋白质，几分钟后 RNA 本身就被销毁，底本安然无恙；而 DNA 复制错一个字母，是会传给子子孙孙的档案事故。两条流水线采用不同的质检标准，恰恰是因为错误的\textbf{代价}不同——又是那本账。

\section{回到那张菜谱}

现在兑现开场的承诺，把你家的祖传菜谱重新讲一遍。

那本书就是 DNA：孤本、珍贵、绝不能冒险带进厨房——正如 DNA 几乎不离开细胞核，被双层核膜护着。你抄在便签纸上的那一页，是 mRNA：内容照抄自底本，但用的是“简装纸”——寿命短、用完即弃，被油污弄脏了就扔，毫无损失。“从此页抄起”的页签，是启动子；执笔的你，是 RNA 聚合酶；菜谱末尾的“完”字，是终止信号。

差别也有几处。你抄菜谱是一次一份，细胞的聚合酶可以在同一个启动子上排队，一个接一个地抄，同一个基因同时有好几份副本在做；你抄完直接能用，真核细胞的副本还要戴帽、加尾、剪接，像便签纸过了塑封、装订才准出核；你把便签带进厨房自己照着做，细胞却把 mRNA 交给核糖体这台通用机器——而核糖体读的那套字母表，和菜谱上的完全不是同一种文字。

\section{便条送达之后}

一个现实应用，先看看这套知识已经怎样改变了医学。你大概听说过 mRNA 疫苗：它的思路直白得漂亮——既然核糖体只认 mRNA，那就人工合成一段 mRNA，编码病毒的一小块特征蛋白，送进人体细胞；你的核糖体照章办事，把这段“外来便条”翻译成蛋白质，免疫系统于是提前见过敌人的长相，学会了识别（第 61 章讲过免疫记忆）。便条本身很快降解，病毒却永远没进过你的身体。没有这一章讲的每个环节——启动子、帽子、尾巴、RNA 的短命特性——这项技术都无从谈起。科学家甚至要给疫苗 mRNA 换上特制的帽子和尾巴，它才活得够久、翻译得够多。

而便条送到核糖体之后，真正的难题才刚开始。mRNA 上还是核酸的语言：A、U、G、C 四个字母排成一列；要造的蛋白质却是二十种氨基酸串成的链。四个字母怎么指挥二十种零件？三个字母一组够吗？谁来当翻译？这正是第 69 章要拆解的问题——翻译：四种字母怎样变成二十种氨基酸。抄好了图纸，接下来该读图纸开工了。

\begin{tiaozhan}
有个细节你可能憋了很久：为什么 DNA 用 T，RNA 却用 U？这两个字母差别极小（T 比 U 多一个甲基），为什么底本和副本非要分开用？一个被广泛接受的解释，妙在它把“短命”和“长寿”的分工用到了极致。胞嘧啶 C 在水里会缓慢地自发脱氨，变成尿嘧啶 U——这是 DNA 档案库最常见的“纸张变质”之一。如果 DNA 本来就用 U 当正常字母，修复酶看见一个 U 将无从判断：这是正常字母，还是 C 变质来的？而 DNA 选择了 T——一个“贴了标签的 U”——之后，规则就变得干净利落：\textbf{DNA 里出现 U，一律是事故，统统换掉}。你的细胞里真有一支专职巡逻队（尿嘧啶糖基化酶），每天在基因组里找 U 拆 U。RNA 为什么不在乎？因为它是消耗品，错一两个字母无所谓，犯不着为每个字母多付一个甲基的成本，也省得维护识别系统。你看，连一个字母的选择，背后都是“存档成本”与“出错代价”的权衡。
\end{tiaozhan}

\section*{练一练}

\begin{sikaoti}
\item 画一幅信息流图：从“细胞需要某种蛋白质”开始，画出 DNA、启动子、RNA 聚合酶、mRNA、核糖体之间的先后关系，并在每条箭头旁注明这一步发生在哪里（核内还是核外）、消耗什么、产出什么。
\item 一段 DNA 模板链的序列是 $\text{3'-ATT GCG TAC-5'}$。请写出转录得到的 mRNA 序列，标清 $5'$ 端和 $3'$ 端；再写出 DNA 的另一条链（编码链）的序列，比较它和 mRNA 的异同。
\item 细菌的 mRNA 平均只能活几分钟，人体细胞的 mRNA 能活几小时。用“底本与副本”的分工逻辑解释：为什么让副本短命反而是优点？如果一个突变让某份 mRNA 永远不会被降解，可能带来什么麻烦？
\item 第 67 章的 DNA 复制和本章的转录都由“聚合酶”完成，工序却大不一样。列一张对照表：模板用量（整份基因组还是局部）、产物份数、是否校对、出错代价，并解释每一行差异背后的原因。
\item 果蝇 Dscam 基因能产生上万种蛋白质版本，但果蝇的基因组并没有因此变大一万倍。用可变剪接的知识，画一幅“同一段原始转录本剪出两种成品”的示意图（用方框表示外显子、直线表示内含子即可）。
\item 逆转录病毒（如 HIV）以 RNA 为模板合成 DNA。据此判断：“遗传信息只能从 DNA 流向 RNA”这句话错在哪？这对你理解“中心法则是主干而非铁律”有什么提醒？
\end{sikaoti}
